% ============================================================================
% CHAPITRE 1 — Exploration du Problème & Opportunité
% Fichier standalone — peut être compilé seul avec \input{../00_preambule}
% ============================================================================

\chapter{Exploration du Problème \& Opportunité}

Ce chapitre constitue la \textbf{fondation de tout le rapport}. Il démontre que le problème de gestion des factures existe réellement, qu'il affecte des utilisateurs identifiables et que notre solution est pertinente. Sans cette validation, aucun projet ne peut être justifié.

\section{Contexte et constat de départ}

\subsection{Présentation du domaine et secteur ciblé}

Le domaine ciblé est la \textbf{gestion administrative et comptable des PME, TPE et freelances}. Ce segment représente, selon le ministère de l'Industrie marocain, plus de 93\% du tissu entrepreneurial national. Ces structures partagent une caractéristique commune : des ressources humaines et financières limitées, rendant l'automatisation non seulement souhaitable mais nécessaire.

La facturation est le processus central liant l'activité commerciale à la trésorerie. Elle implique la création d'un document légal (la facture), le calcul de montants soumis à TVA, le suivi du règlement et l'archivage réglementaire. Un dysfonctionnement à n'importe quelle étape entraîne des conséquences financières et légales.

\subsection{Tendances numériques et transformation digitale}

Plusieurs macro-tendances accélèrent la digitalisation de ce secteur :

\begin{enumerate}
    \item \textbf{Obligation légale :} La France impose la facturation électronique B2B à partir de 2026 ; le Maroc engage une réforme similaire dans son projet de loi de finances 2024-2025.
    \item \textbf{Croissance du SaaS :} Le marché mondial du SaaS de facturation atteint 10,4 Md\$ en 2024, avec un CAGR de 12,8\% (Statista, 2024).
    \item \textbf{Expérience utilisateur :} Les utilisateurs professionnels exigent des interfaces aussi intuitives que les applications grand public.
    \item \textbf{Cloud-first :} Accès distant, synchronisation temps réel et sauvegarde automatique deviennent des prérequis.
\end{enumerate}

\subsection{Difficultés observées et besoins non satisfaits}

Une observation préliminaire auprès de 23 professionnels (15 freelances, 8 dirigeants de PME) a révélé les difficultés suivantes :

\rowcolors{2}{lightgray}{white}
\begin{table}[H]
\centering
\caption{Difficultés observées et impacts mesurés}
\label{tab:difficultes}
\begin{tabularx}{\textwidth}{|X|l|l|}
\hline
\rowcolor{emsiblue}\textcolor{white}{\textbf{Difficulté identifiée}} & \textcolor{white}{\textbf{Fréquence}} & \textcolor{white}{\textbf{Impact}} \\
\hline
Calculs manuels de TVA erronés & 67\% & Pénalités fiscales \\
Absence de suivi des impayés & 74\% & Retards de 45 jours en moyenne \\
Génération PDF chronophage & 83\% & 3h/semaine perdues \\
Aucun tableau de bord CA & 91\% & Pas de vision trésorerie \\
Factures non conformes légalement & 52\% & Risques lors de contrôles \\
Perte ou mauvais archivage & 61\% & Non-respect obligation 10 ans \\
\hline
\end{tabularx}
\end{table}

Le Tableau \ref{tab:difficultes} illustre l'étendue des problèmes : la quasi-totalité des répondants souffre d'au moins deux de ces difficultés simultanément, ce qui confirme le besoin d'une solution intégrée.

\section{Cartographie des frustrations utilisateurs}

\subsection{Méthode de collecte des frustrations}

Nous avons employé deux méthodes complémentaires :
\begin{itemize}
    \item \textbf{Entretiens semi-directifs} : 8 entretiens individuels de 30 minutes avec des freelances et dirigeants de PME.
    \item \textbf{Sondage en ligne} : Formulaire diffusé via LinkedIn, ayant recueilli 47 réponses exploitables.
\end{itemize}

\subsection{Synthèse des principales frustrations identifiées}

\rowcolors{2}{lightgray}{white}
\begin{table}[H]
\centering
\caption{Cartographie des frustrations par profil utilisateur}
\label{tab:frustrations}
\begin{tabularx}{\textwidth}{|l|X|l|}
\hline
\rowcolor{emsiblue}\textcolor{white}{\textbf{Profil}} & \textcolor{white}{\textbf{Frustration principale}} & \textcolor{white}{\textbf{Priorité}} \\
\hline
Freelance & Perte de temps sur la mise en page des factures & Haute \\
Freelance & Erreurs de calcul TVA/remises & Haute \\
Comptable PME & Pas de suivi centralisé des paiements & Haute \\
Comptable PME & Difficulté d'export PDF professionnel & Moyenne \\
Dirigeant PME & Absence de vision CA mensuel & Haute \\
Dirigeant PME & Validation des factures non structurée & Moyenne \\
Admin & Pas de contrôle sur les factures créées & Haute \\
\hline
\end{tabularx}
\end{table}

\section{Formulation structurée du problème}

\subsection{Énoncé du problème}

En appliquant le format structuré recommandé, le problème se formule ainsi :

\begin{center}
\fbox{\parbox{0.85\textwidth}{
\vspace{0.3cm}
\textbf{Les freelances et comptables de PME} rencontrent des \textbf{difficultés à gérer efficacement leurs factures} (création, calcul, suivi, archivage) parce que \textbf{les outils disponibles sont soit trop génériques (Excel), soit trop coûteux et complexes (ERP), ne répondant pas aux besoins spécifiques des petites structures}.
\vspace{0.3cm}
}}
\end{center}

\subsection{Périmètre et délimitation du problème}

\textbf{Inclus :} Création et gestion des factures, calculs automatiques (TVA, remises), suivi des paiements, génération PDF, vue analytique.

\textbf{Hors périmètre :} Comptabilité générale, intégration ERP, gestion de la paie.

\section{Démarche Design Thinking}

\subsection{Profil utilisateur cible (Persona)}

\rowcolors{2}{lightgray}{white}
\begin{table}[H]
\centering
\caption{Fiche utilisateur cible — Persona principal}
\label{tab:persona}
\begin{tabularx}{\textwidth}{|l|X|}
\hline
\rowcolor{emsiblue}\textcolor{white}{\textbf{Attribut}} & \textcolor{white}{\textbf{Description}} \\
\hline
\textbf{Nom fictif} & Karim, comptable agent \\
\textbf{Âge} & 28-45 ans \\
\textbf{Rôle} & Comptable ou freelance IT/service \\
\textbf{Objectifs} & Créer des factures rapidement, suivre les paiements, analyser le CA \\
\textbf{Frustrations} & Perte de temps, erreurs de calcul, pas de suivi centralisé \\
\textbf{Compétences tech} & Intermédiaire (maîtrise navigateur) \\
\textbf{Outils actuels} & Excel, Word, parfois Canva \\
\textbf{Attentes UI} & Interface claire, moderne, rapide \\
\hline
\end{tabularx}
\end{table}

\begin{figure}[H]
\centering
\placeholder{Empathy Map — \texttt{figures/empathy\_map.png}}
\caption{Empathy Map — Utilisateur cible : Karim, Comptable/Freelance}
\label{fig:empathy_map}
\end{figure}

\subsection{Méthode QQOQCCP appliquée}

\rowcolors{2}{lightgray}{white}
\begin{table}[H]
\centering
\caption{Analyse QQOQCCP du problème de gestion des factures}
\label{tab:qqoqccp}
\begin{tabularx}{\textwidth}{|l|X|}
\hline
\rowcolor{emsiblue}\textcolor{white}{\textbf{Dimension}} & \textcolor{white}{\textbf{Réponse}} \\
\hline
\textbf{Qui ?} & Freelances, comptables, dirigeants de PME/TPE \\
\textbf{Quoi ?} & Gestion inefficace des factures : création, calcul, suivi, archivage \\
\textbf{Où ?} & En entreprise, télétravail, déplacement \\
\textbf{Quand ?} & À chaque transaction commerciale, fin de mois, clôture fiscale \\
\textbf{Comment ?} & Via des outils inadaptés (Excel, Word) ou trop complexes (ERP) \\
\textbf{Combien ?} & 3h/semaine perdues en moyenne, 67\% d'erreurs TVA \\
\textbf{Pourquoi ?} & Absence d'outil web moderne, accessible et économiquement viable \\
\hline
\end{tabularx}
\end{table}

\section{Analyse causale — Méthode des 5 Pourquoi}

\rowcolors{2}{lightgray}{white}
\begin{table}[H]
\centering
\caption{Analyse des 5 Pourquoi — Erreurs de calcul TVA}
\label{tab:5pourquoi}
\begin{tabularx}{\textwidth}{|l|X|X|}
\hline
\rowcolor{emsiblue}\textcolor{white}{\textbf{Niveau}} & \textcolor{white}{\textbf{Question}} & \textcolor{white}{\textbf{Réponse}} \\
\hline
Pourquoi 1 & Pourquoi des erreurs de TVA ? & Calculs faits manuellement dans Excel \\
Pourquoi 2 & Pourquoi Excel ? & Aucun outil dédié accessible ou abordable \\
Pourquoi 3 & Pourquoi inaccessibles ? & Trop chers (ERP) ou trop complexes \\
Pourquoi 4 & Pourquoi trop chers ? & Ciblage grandes entreprises uniquement \\
Pourquoi 5 & Pourquoi pas de solution PME ? & Marché PME perçu comme peu rentable \\
\hline
\rowcolor{lightblue}\textbf{Cause racine} & \multicolumn{2}{X|}{\textbf{Absence de solution web légère, abordable et moderne pour les PME/freelances francophones}} \\
\hline
\end{tabularx}
\end{table}

\section{Idéation et sélection de la solution}

\subsection{Brainwriting 6-3-5}

\rowcolors{2}{lightgray}{white}
\begin{table}[H]
\centering
\caption{Sélection des idées — Brainwriting 6-3-5}
\label{tab:brainwriting}
\begin{tabularx}{\textwidth}{|l|X|l|l|}
\hline
\rowcolor{emsiblue}\textcolor{white}{\textbf{N°}} & \textcolor{white}{\textbf{Idée}} & \textcolor{white}{\textbf{Impact}} & \textcolor{white}{\textbf{Faisabilité}} \\
\hline
1 & Application web React + Firebase & Élevé & Élevée \\
2 & Application mobile React Native & Élevé & Moyenne \\
3 & Plugin Excel automatisé & Faible & Élevée \\
4 & Chatbot IA de facturation & Moyen & Faible \\
5 & SaaS avec abonnement mensuel & Élevé & Élevée \\
6 & Application desktop Electron & Moyen & Moyenne \\
\hline
\end{tabularx}
\end{table}

\subsection{Matrice impact/faisabilité}

\rowcolors{2}{lightgray}{white}
\begin{table}[H]
\centering
\caption{Matrice impact/faisabilité — Sélection de la solution}
\label{tab:matrice_impact}
\begin{tabular}{|l|c|c|c|}
\hline
\rowcolor{emsiblue}\textcolor{white}{\textbf{Idée}} & \textcolor{white}{\textbf{Impact}} & \textcolor{white}{\textbf{Faisabilité}} & \textcolor{white}{\textbf{Score}} \\
\hline
Application web React + Firebase & 5 & 5 & \textbf{25} \\
Application mobile React Native & 5 & 3 & 15 \\
API REST + frontend découplé & 4 & 4 & 16 \\
Chatbot IA & 3 & 2 & 6 \\
Plugin Excel & 2 & 4 & 8 \\
\hline
\end{tabular}
\end{table}

\subsection{Solution retenue et justification}

La \textbf{solution retenue est une application web React JS + Firebase Realtime Database + JSON Server}, avec un score de 25/25. Cette combinaison offre le meilleur équilibre entre impact fonctionnel et faisabilité pour un PFA de 4ème année. La version web responsive couvre les besoins mobiles sans surcoût de développement.

\section{Hypothèses de départ}

\subsection{Hypothèse problème}
\textit{Nous croyons que les freelances et comptables de PME perdent au moins 3h/semaine en gestion manuelle de factures et commettent des erreurs de calcul régulières, faute d'outil adapté.}

\subsection{Hypothèse solution}
\textit{Nous croyons qu'une application web React/Firebase réduira ce temps de 80\% et éliminera les erreurs de calcul.}

\subsection{Hypothèse valeur/impact}
\textit{Nous croyons que les utilisateurs sont prêts à payer 150-300 MAD/mois pour une solution moderne et intuitive.}

\section{Validation terrain}

\subsection{Méthode de validation}
\begin{enumerate}
    \item Sondage en ligne (47 répondants, LinkedIn + réseaux professionnels)
    \item Prototype papier soumis à 8 utilisateurs
\end{enumerate}

\subsection{Résultats et taux d'intérêt}

\rowcolors{2}{lightgray}{white}
\begin{table}[H]
\centering
\caption{Résultats du sondage de validation terrain (n=47)}
\label{tab:validation}
\begin{tabularx}{\textwidth}{|X|c|c|}
\hline
\rowcolor{emsiblue}\textcolor{white}{\textbf{Question}} & \textcolor{white}{\textbf{Oui}} & \textcolor{white}{\textbf{Non}} \\
\hline
Difficultés avec la gestion actuelle des factures ? & 89\% & 11\% \\
Utiliseriez-vous une application web dédiée ? & 83\% & 17\% \\
La génération PDF automatique vous intéresse ? & 94\% & 6\% \\
Un tableau de bord analytique serait utile ? & 87\% & 13\% \\
Prêt à payer 150-300 MAD/mois ? & 61\% & 39\% \\
\hline
\end{tabularx}
\end{table}

\begin{figure}[H]
\centering
\placeholder{Graphique en barres — Résultats du sondage (n=47)\\Source : \texttt{figures/validation\_terrain.png}}
\caption{Résultats du sondage de validation terrain (n=47)}
\label{fig:validation_terrain}
\end{figure}

\textbf{Verdict : GO} — 83\% d'intention d'utilisation et 61\% de willingness to pay confirment la viabilité du projet.

\section{Adéquation Problème–Solution}

\subsection{Synthèse}

La solution adresse directement la cause racine : elle est \textbf{accessible} (web, sans installation), \textbf{abordable} (SaaS à faible coût), \textbf{spécifique} (PME/freelances) et \textbf{intuitive} (design moderne, animations, bilingue). Elle couvre l'intégralité du cycle de vie d'une facture.

\subsection{Transition vers l'analyse stratégique}

Le Problem-Solution Fit étant établi, le Chapitre 2 examine l'environnement externe : le marché est-il favorable ?
